% =====================================================================
%  ANEXOS
%  Responsables: INTEGRANTE 2 (Anexos A y E) e INTEGRANTE 4 (Anexos B, C y D)
%  Extensión: libre, no cuenta para el objetivo de páginas del cuerpo.
%
%  DOS COSAS IMPORTANTES DE ESTE FORMATO:
%
%  1. main.tex conmuta a UNA COLUMNA (\onecolumn) antes de incluir este
%     archivo. Por eso aquí SÍ se puede usar longtable, que no funciona
%     en modo de dos columnas. Es el único sitio del informe donde se
%     puede.
%
%  2. Los anexos se rotulan a mano (Anexo A, B, C...) con \chapter* y se
%     añaden al índice con \addcontentsline. No se usa \appendix: en
%     IEEEtran, combinado con el adaptador de jerarquía del preámbulo,
%     produce resultados impredecibles.
% =====================================================================

\newpage
\chapter*{Anexo A. Mapa de ubicación ampliado}
\addcontentsline{toc}{section}{Anexo A. Mapa de ubicación ampliado}
\label{anexo:mapa}

% TODO: Versión a página completa del mapa de la Figura 1, con el detalle de
%   los municipios de San Miguel Ixtahuacán y Sipacapa, las cuencas de los
%   ríos Tzalá, Quivichil y Cuilco, y la ubicación de la planta y del
%   depósito de relaves.
%   RECORDATORIO DE LICENCIA: si el mapa base es de OpenStreetMap, la
%   atribución «© OpenStreetMap contributors» (licencia ODbL) es obligatoria.
%   Si usas el mapa del conflicto de OCMAL, cítalo. \cite{ocmalsfmarlin}
%
%\begin{figure}
%  \centering
%  \includegraphics[width=0.9\textwidth]{imagenes/mapas/mapa_ampliado.png}
%  \caption{Mapa ampliado del área de influencia de la Mina Marlin}
%  \label{fig:mapa-ampliado}
%\end{figure}

\newpage
\chapter*{Anexo B. Matriz de no conformidades completa}
\addcontentsline{toc}{section}{Anexo B. Matriz de no conformidades completa}
\label{anexo:matriz}

% TODO: En el capítulo 6 va un EXTRACTO de la matriz (los hallazgos
%   principales, en una tabla que abarca las dos columnas). Aquí va la
%   matriz ÍNTEGRA, con todos los hallazgos de los siete ejes.
%   Como estamos a una columna, aquí sí funciona longtable.
%
%\begin{landscape}
%\begin{footnotesize}
%\begin{longtable}{@{}L{1.8cm} L{4cm} L{4cm} L{3cm} L{1.8cm} L{4cm} L{2cm}@{}}
%  \caption{Matriz completa de no conformidades de la auditoría socioambiental de la Mina Marlin}
%  \label{tab:matriz-nc-completa} \\
%  \toprule
%  \textbf{Eje} & \textbf{Hallazgo} & \textbf{Evidencia objetiva} & \textbf{Criterio incumplido} & \textbf{Clasif.} & \textbf{Recomendación} & \textbf{Estado de cierre} \\
%  \midrule
%  \endfirsthead
%  \multicolumn{7}{l}{\footnotesize\textit{Tabla \ref{tab:matriz-nc-completa} (continuación)}} \\
%  \toprule
%  \textbf{Eje} & \textbf{Hallazgo} & \textbf{Evidencia objetiva} & \textbf{Criterio incumplido} & \textbf{Clasif.} & \textbf{Recomendación} & \textbf{Estado de cierre} \\
%  \midrule
%  \endhead
%  \bottomrule
%  \multicolumn{7}{l}{\footnotesize Fuente: elaboración propia a partir de \cite{oncommonground2010hria}.} \\
%  \endlastfoot
%  Consulta    & & & Convenio 169 {OIT}, art. 6 & \NCmayor & & \\
%  Ambiente    & & &                            &          & & \\
%  Trabajo     & & &                            &          & & \\
%  Tierras     & & & {IFC} {PS}5                &          & & \\
%  Inversión   & & &                            &          & & \\
%  Seguridad   & & & {VPSHR}                    &          & & \\
%  Remediación & & & {IFC} {PS}1                &          & & \\
%\end{longtable}
%\end{footnotesize}
%\end{landscape}

\newpage
\chapter*{Anexo C. Cronología del caso (1996--2017)}
\addcontentsline{toc}{section}{Anexo C. Cronología del caso (1996--2017)}
\label{anexo:cronologia}

% TODO: Tabla de dos columnas (fecha / hecho). Hitos que no pueden faltar:
%   06/08/1996  Constitución de Montana Exploradora de Guatemala, S.A.
%   1998        Descubrimiento del yacimiento.
%   2000        Adquisición por Francisco Gold.
%   2002        Fusión con Glamis Gold.
%   2003        Préstamo de la IFC por US$ 45 millones.
%   27/11/2003  Licencia de explotación del Ministerio de Energía y Minas.
%   2004-2005   Construcción de la mina.
%   2004        Revisión independiente del EIA por Robert Moran.
%   2005        Consulta comunitaria de Sipacapa; queja de Madre Selva ante la CAO.
%   Fines 2005  Inicio de producción.
%   05/2006     Cierre del caso CAO Marlin-01/Sipacapa.
%   2006        Goldcorp adquiere Glamis Gold.
%   03/2008     Firma del MOU para la evaluación de derechos humanos.
%   10/2008     Selección de On Common Ground Consultants.
%   05/2010     Publicación de la HRIA.
%   20/05/2010  La CIDH otorga medidas cautelares MC 260-07.
%   2010        Estudio de salud de Physicians for Human Rights / U. de Michigan.
%   2010        Pronunciamiento de la CEACR de la OIT.
%   12/2011     La CIDH modifica las medidas cautelares.
%   2013        Instalación de la planta de relaves filtrados (dry-stack).
%   31/05/2017  Cierre de producción de la mina.
%
%\begin{longtable}{@{}L{3cm} L{12cm}@{}}
%  \caption{Cronología del caso Mina Marlin} \label{tab:cronologia} \\
%  \toprule \textbf{Fecha} & \textbf{Hecho} \\ \midrule \endfirsthead
%  \toprule \textbf{Fecha} & \textbf{Hecho} \\ \midrule \endhead
%  \bottomrule \endlastfoot
%   & \\
%\end{longtable}

\newpage
\chapter*{Anexo D. Fichas de los estándares aplicados}
\addcontentsline{toc}{section}{Anexo D. Fichas de los estándares aplicados}
\label{anexo:estandares}

% TODO: Una ficha breve por estándar, con: nombre completo, organismo emisor,
%   año, carácter (obligatorio / voluntario / certificable), alcance y
%   pertinencia para este caso. Estándares a fichar:
%   ISO 14001:2015, ISO 19011:2018, IFC Performance Standards,
%   Convenio 169 de la OIT, VPSHR, ICMM Mining Principles, GISTM e IRMA.
%   Sugerencia: la macro \recuadro{título}{contenido} del preámbulo da un
%   formato de ficha limpio y homogéneo.

\newpage
\chapter*{Anexo E. Glosario de siglas}
\addcontentsline{toc}{section}{Anexo E. Glosario de siglas}
\label{anexo:glosario}

% TODO: Lista alfabética. Como mínimo:
%   AMAC    Asociación de Monitoreo Ambiental Comunitario
%   ARD     Drenaje ácido de roca (acid rock drainage)
%   CAO     Compliance Advisor Ombudsman
%   CEACR   Comisión de Expertos en Aplicación de Convenios y Recomendaciones (OIT)
%   CIDH    Comisión Interamericana de Derechos Humanos
%   CIL     Carbon in leach (lixiviación con carbón en pulpa)
%   COPAE   Comisión Pastoral Paz y Ecología
%   DIHR    Danish Institute for Human Rights
%   EIA     Estudio de impacto ambiental
%   GISTM   Global Industry Standard on Tailings Management
%   HRCA    Human Rights Compliance Assessment
%   HRIA    Human Rights Impact Assessment
%   ICMM    International Council on Mining and Metals
%   IFC     International Finance Corporation
%   IRMA    Initiative for Responsible Mining Assurance
%   MOU     Memorando de entendimiento
%   NC      No conformidad
%   OCMAL   Observatorio de Conflictos Mineros de América Latina
%   OEFA    Organismo de Evaluación y Fiscalización Ambiental (Perú)
%   OIT     Organización Internacional del Trabajo
%   PAMA    Programa de adecuación y manejo ambiental
%   PHVA    Planificar, Hacer, Verificar, Actuar
%   PS      Performance Standard (IFC)
%   SEIA    Sistema Nacional de Evaluación de Impacto Ambiental
%   SGA     Sistema de gestión ambiental
%   TSF     Tailings storage facility (depósito de relaves)
%   VPSHR   Voluntary Principles on Security and Human Rights

\pendiente{completar los anexos A a E}
